% !TeX program = pdflatex
% !BIB program = bibtex

\documentclass[conference]{IEEEtran}
\IEEEoverridecommandlockouts

\usepackage{amsmath,amssymb,amsfonts}
\usepackage{booktabs}
\usepackage{cite}
\usepackage{graphicx}
\usepackage{textcomp}
\usepackage{xcolor}

\def\BibTeX{{\rm B\kern-.05em{\sc i\kern-.025em b}\kern-.08em
    T\kern-.1667em\lower.7ex\hbox{E}\kern-.125emX}}

\begin{document}

\title{Two-Timescale Fair Downlink Resource Allocation with Handover Budgets in LEO Satellite Networks}

\author{
\IEEEauthorblockN{First Author, Second Author, and Third Author}
\IEEEauthorblockA{
\textit{Department or Laboratory}\\
\textit{University or Organization}\\
City, Country\\
email@example.com}
}

\maketitle

\begin{abstract}
Low Earth orbit (LEO) direct-access systems must allocate scarce downlink
resources while managing frequent satellite-cell association changes. Existing
work often optimizes radio resources or handovers separately, or handles
handover through soft penalties. This paper studies a two-timescale optimization
framework for fair LEO downlink service under hard per-cell handover budgets.
At the fast timescale, physical resource blocks and transmit power are jointly
allocated for a fixed satellite-cell association through a convex proportional
fairness problem. At the slow timescale, satellite-cell associations are
selected over a visibility horizon while respecting explicit handover budgets.
The resulting framework supports exact small-scale optimization and scalable
rolling-horizon approximation. The evaluation is designed to quantify served
demand, proportional fairness, worst-cell satisfaction, handover compliance, and
solver runtime.
\end{abstract}

\begin{IEEEkeywords}
LEO satellite networks, resource allocation, handover budget, convex
optimization, mixed-integer programming, proportional fairness
\end{IEEEkeywords}

% !TeX root = ../main.tex

\section{Introduction}

Low Earth orbit (LEO) satellite networks are a promising architecture for
wide-area direct-access connectivity, but their downlink scheduling problem is
shaped by two coupled sources of variation. At the fast timescale, cells compete
for physical resource blocks (PRBs) and satellite transmit power under changing
demand. At the slow timescale, satellite visibility evolves with orbital motion,
which forces satellite-cell association decisions and may trigger frequent
handovers.

Existing LEO downlink allocation studies have emphasized fair resource
balancing across large service regions~\cite{leyva2023continent}, while
handover-aware formulations often use soft switching penalties
\cite{afif2025handover}. These approaches are valuable baselines, but two
practical gaps remain. First, PRB and power allocation should be optimized
jointly when evaluating fixed associations. Second, a soft handover penalty does
not guarantee that each cell stays within a deployment-level handover budget.
This motivates a two-timescale formulation that combines fast convex resource
allocation with slow hard-budget association optimization.

This paper focuses on the optimization core and leaves online learning under
uncertain demand as future work. The intended contributions are:
\begin{itemize}
\item a fast-timescale convex PRB--power allocation model for fixed LEO
associations;
\item a slow-timescale satellite-cell association formulation with hard
per-cell handover budgets;
\item exact and rolling-horizon solution methods that trade optimality for
runtime;
\item a reproducible simulation plan measuring throughput, proportional
fairness, worst-cell satisfaction, handover count, and solve time.
\end{itemize}

% !TeX root = ../main.tex

\section{System Model}

Consider a LEO constellation with satellite set $\mathcal{S}$ and cell set
$\mathcal{C}$. Time is divided into slow slots $k \in \mathcal{K}$, and each
slow slot contains fast slots $m \in \mathcal{M}$. Binary visibility
$v_{s,c,k}$ indicates whether satellite $s$ can serve cell $c$ during slow slot
$k$. The slow-timescale decision is the binary association
$x_{s,c,k} \in \{0,1\}$, with each cell assigned to exactly one visible
satellite.

Each satellite $s$ has PRB budget $N_s$ and transmit-power budget $P_s$. During
slow slot $k$, cell $c$ receives demand arrivals $a_{c,k,m}$ over fast slots.
For a fixed slow-slot association, the fast allocator chooses PRB allocation
$n_{s,c,k,m}$ and power allocation $p_{s,c,k,m}$. With channel gain
$g_{s,c,k}$, the served demand is limited by the Shannon-type downlink rate
\begin{equation}
    r_{s,c,k,m}
    =
    n_{s,c,k,m} W
    \log_2 \left(1 +
    \frac{g_{s,c,k} p_{s,c,k,m}}
    {n_{s,c,k,m} N_0 W}
    \right),
\end{equation}
where $W$ is the bandwidth per PRB and $N_0$ is the noise spectral density.

Let $z_{c,k,m}$ denote the served demand of cell $c$ in fast slot $m$. The
slow-slot satisfaction rate is
\begin{equation}
    \xi_{c,k}
    =
    \frac{\sum_{m \in \mathcal{M}} z_{c,k,m}}
    {\sum_{m \in \mathcal{M}} a_{c,k,m} + \epsilon_a},
\end{equation}
where $\epsilon_a$ avoids division by zero. Proportional fairness is represented
by the log utility
\begin{equation}
    U = \sum_{c \in \mathcal{C}} \sum_{k \in \mathcal{K}}
    \log(\epsilon + \xi_{c,k}),
\end{equation}
where $\epsilon > 0$ stabilizes zero-service cases.

Handovers are counted per cell. Let $h_{c,k}$ indicate whether cell $c$ changes
its serving satellite between slow slots $k-1$ and $k$. The deployment
constraint is a hard budget:
\begin{equation}
    \sum_{k \in \mathcal{K}} h_{c,k} \leq H_c,\quad \forall c \in \mathcal{C}.
\end{equation}

% !TeX root = ../main.tex

\section{Problem Formulation}

The full problem couples binary association decisions with continuous PRB and
power allocation. Directly solving this mixed-integer nonlinear program at paper
scale is not practical, so we expose its structure through a two-timescale
decomposition.

\subsection{Fast Allocation for Fixed Association}

Given $x_{s,c,k}$, the fast-timescale problem optimizes PRB and power variables
for each slow slot. Its objective is proportional fairness:
\begin{equation}
    \max_{\{n,p,z,\xi\}}
    \sum_{c \in \mathcal{C}} \log(\epsilon + \xi_{c,k}).
\end{equation}
The constraints are
\begin{align}
    &\sum_{c \in \mathcal{C}} n_{s,c,k,m} \leq N_s,
    &&\forall s,m,\\
    &\sum_{c \in \mathcal{C}} p_{s,c,k,m} \leq P_s,
    &&\forall s,m,\\
    &n_{s,c,k,m}=p_{s,c,k,m}=0,
    &&\text{if } x_{s,c,k}=0,\\
    &0 \leq z_{c,k,m} \leq a_{c,k,m},
    &&\forall c,m.
\end{align}
The rate expression is handled in perspective form, which makes the rate
constraint concave in $(n,p)$ and yields a disciplined convex program for fixed
association.

\subsection{Slow Association with Handover Budgets}

At the slow timescale, the optimizer selects $x_{s,c,k}$ subject to visibility,
single-association, and handover constraints:
\begin{align}
    &\sum_{s \in \mathcal{S}} x_{s,c,k}=1,
    &&\forall c,k,\\
    &x_{s,c,k}\leq v_{s,c,k},
    &&\forall s,c,k,\\
    &h_{c,k}\geq x_{s,c,k}-x_{s,c,k-1},
    &&\forall s,c,k>1,\\
    &h_{c,k}\geq x_{s,c,k-1}-x_{s,c,k},
    &&\forall s,c,k>1,\\
    &\sum_{k>1} h_{c,k}\leq H_c,
    &&\forall c.
\end{align}
The association objective uses a calibrated utility proxy
$q_{s,c,k}$, obtained from channel quality, expected demand load, and the fast
allocation kernel. This gives the slow problem
\begin{equation}
    \max_{\{x,h\}}
    \sum_{c,k,s} q_{s,c,k} x_{s,c,k}
    - \lambda_h \sum_{c,k} h_{c,k},
\end{equation}
while retaining the hard budget constraint for every cell.

% !TeX root = ../main.tex

\section{Proposed Method}

\subsection{Inner Convex Resource Allocation}

For every candidate association or calibration sample, the inner allocator
solves the fixed-association PRB--power problem. The key modeling step is the
perspective transformation of the rate term. In natural logarithm form,
\begin{equation}
    n \log \left(1 + \alpha p/n\right)
    =
    -\mathrm{rel\_entr}(n,n+\alpha p),
\end{equation}
which is concave in $(n,p)$. This enables a convex ground-truth evaluator for
service satisfaction under fixed association.

\subsection{Outer Handover-Constrained Optimization}

The outer layer optimizes the association trajectory using the calibrated proxy
$q_{s,c,k}$. For small and medium instances, we solve the binary problem as a
mixed-integer program to obtain an offline benchmark. For larger horizons, we
use a rolling-horizon variant: solve a shorter window, commit the leading
slots, carry forward the last association, and update the remaining handover
budget. This preserves hard handover feasibility while reducing solve time.

\subsection{Algorithm Summary}

The complete workflow is:
\begin{enumerate}
\item generate visibility, channel, demand, and handover-budget inputs;
\item estimate $q_{s,c,k}$ using the inner PRB--power allocation kernel or a
calibrated capacity proxy;
\item solve the slow association problem exactly or with rolling horizon;
\item evaluate the selected association with the same fast allocation kernel;
\item report fairness, served demand, handovers, budget violations, and runtime.
\end{enumerate}

% !TeX root = ../main.tex

\section{Simulation Setup and Results}

The evaluation will compare all methods on shared LEO downlink scenario
instances. The planned baselines are nearest-visible-satellite association,
best-channel association without handover control, soft handover-penalized
association, full-horizon mixed-integer association, and rolling-horizon
association. When the proxy is separable across cells, a dynamic-programming
solver can also be included as an exact reference for the surrogate objective.

\begin{table}[t]
\caption{Primary Evaluation Metrics}
\label{tab:metrics}
\centering
\begin{tabular}{ll}
\toprule
Metric & Purpose \\
\midrule
Mean satisfaction & Average served-demand ratio \\
Fifth-percentile satisfaction & Worst-cell service quality \\
Log utility & Proportional fairness \\
Total served demand & System throughput \\
Total handovers & Mobility overhead \\
Budget violations & Hard-constraint feasibility \\
Solve time & Algorithm scalability \\
\bottomrule
\end{tabular}
\end{table}

The main metrics are summarized in Table~\ref{tab:metrics}. The final WCNC
claim should be limited to scenarios where the proposed method has zero
avoidable handover-budget violations. Sensitivity studies should vary:
\begin{itemize}
\item demand scale and demand imbalance across cells;
\item visibility density and minimum elevation threshold;
\item handover budget $H_c$;
\item rolling-window length and commit step;
\item proxy calibration error relative to the inner convex allocator.
\end{itemize}

% !TeX root = ../main.tex

\section{Conclusion}

This paper targets handover-constrained fair downlink allocation in LEO
satellite networks. The proposed two-timescale structure combines a convex
fast-timescale PRB--power allocator with a slow-timescale association optimizer
that enforces hard per-cell handover budgets. The planned evaluation focuses on
served demand, proportional fairness, worst-cell satisfaction, handover
compliance, and runtime. Online scheduling under demand uncertainty and
hierarchical reinforcement learning are natural extensions, but are kept outside
the core WCNC submission scope.


\bibliographystyle{IEEEtran}
\bibliography{../bib/leo_references}

\end{document}
